\section{Artifact and Scale}
\label{sec:artifact-scale}

The project is a large Lean artifact, not a short isolated formalization.  The
active \code{Stokes/} directory contains 544 Lean files and approximately
165006 Lean lines.  This scale reflects both the mathematics and the proof
engineering required to compress a certificate-heavy development into a small
public theorem.  The maintenance and interface issues are in the same family
as those discussed for large mathematical libraries such as mathlib
\cite{vandoorn2020maintaining,baanen2022instance,baanen2025growing}.

\subsection{Repository organization}

The artifact repository is available at
\url{https://github.com/evaristebernhard/Stokes} (accessed 29 May 2026).
The active Lean files live under \code{Stokes/}.  The main public route is
exported by
\[
\code{Stokes.Global.CoverIndexedFirstPrinciples}
\]
and by the slimmed cover-indexed public entry
\[
\code{Stokes.Global.CoverIndexed}.
\]
Historical routes remain available as implementation or compatibility modules,
but the default public surface is centered on the first-principles theorem.

The proof archaeology documents used to prepare this paper live under
\code{docs/formalization-log/}.  They record the mathematical bottlenecks
behind support, zero-localization, half-space boundaries, compact finite
selection, endpoint reconstruction, and API slimming.  These documents are not
part of the trusted proof base, but they are useful for explaining the
artifact.

\subsection{Public API slimming}

During development, several intermediate theorem shapes exposed internal
certificates: raw selected covers, clean inputs, natural inputs, collar routes,
endpoint adapters, image-control packages, local Stokes fields, and manual
support assumptions.  These stages were useful because they stabilized the
statement while independent local obligations were being proved.  The final
API is the result of removing those theorem-facing fields.

The first-principles theorem is therefore best viewed as a compression result,
but not in the superficial sense of putting a small facade around a large
record.  It proves that the records which earlier routes consumed are
generated by the mathematical hypotheses.  In this sense the artifact moves
from a certificate-consuming interface to a certificate-generating interface.
It does not bypass the hidden proof objects.  It constructs them internally and
then applies the intrinsic represented theorem.  This explains why a short
top-level proof can rest on a large body of Lean code.

\subsection{Route evolution}

The development history is visible in route names such as Raw, Clean, Natural,
Intrinsic, and FirstPrinciples.  These names should not be read as competing
theorems of equal status.  They record a field-elimination process.

\begin{center}
\begin{tabular}{ll}
\hline
Route style & Main theorem-facing burden \\
\hline
Raw & selected cover and low-level local data are explicit \\
Clean & several geometric certificates are packaged but still visible \\
Natural & generated data starts replacing manual endpoint fields \\
Intrinsic & pointwise chart-box data, compactness, and support control remain \\
FirstPrinciples & only chartwise smoothness and compact closed support remain \\
\hline
\end{tabular}
\end{center}

The engineering point is that intermediate routes made hidden obligations
visible.  Once an obligation could be generated from earlier mathematical
data, it moved out of the public input.  The final theorem is therefore not
another layer of ceremony; it is the accumulated result of proving that the
former certificate fields are consequences of the natural hypotheses.

\subsection{Verification gates}

The development is checked by the standard Lean build commands.  For the paper
claim, the relevant commands are:
\begin{lstlisting}
lake build Stokes
rg "\bsorry\b|\badmit\b|^\s*axiom\b" --glob "*.lean"
\end{lstlisting}
The theorem-level axiom audit is:
\begin{lstlisting}
#print axioms
  Stokes.CoverIndexedZeroCompactRepresentedStokesFirstPrinciplesInput.representedStokes
\end{lstlisting}
The expected output contains the standard Lean/mathlib axioms:
\[
\code{propext},\quad \code{Classical.choice},\quad \code{Quot.sound}.
\]

This paper rewrite does not rerun the full Lean build, because it modifies only
\code{paper/}.  The Lean names quoted in the text were checked by lightweight
repository search against the active sources.

\subsection{What scale bought}

The artifact's size bought several mathematical reductions:
\begin{itemize}
  \item ordinary total chart support assumptions were replaced by
  zero-localized support theorems;
  \item boundary-chart ambient-open assumptions were replaced by half-space
  interior and closed-box regularity fields;
  \item user-supplied finite covers and refinements were replaced by
  compactness-driven selected constructions;
  \item endpoint adapters were replaced by generated localized refined
  endpoints;
  \item theorem-facing geometric certificates were internalized behind the
  first-principles route.
\end{itemize}

These reductions are the main proof-engineering achievement of the project.
They also make the theorem usable as a stable milestone: future work can build
native integral comparison theorems against the canonical represented
endpoints rather than against a moving collection of intermediate certificates.

The zero-localized support reduction is representative of the whole pattern.
The old theorem-facing hypothesis would have asked the user to control the
support of an ordinary totalized chart representative.  The final route proves
instead that the generated zero-localized representative has the required
topological support.  This removes a bad hypothesis by replacing the object of
the theorem and then proving the comparison and support facts needed to use it.

\subsection{Artifact reading guide}

A reader who wants to audit the mathematics should start from the
first-principles theorem, then follow the intrinsic route backward: canonical
endpoint, intrinsic route certificate, selected smooth refinement, zero support,
interior fields, and half-space local Stokes.  A reader who wants to audit the
engineering should instead compare the historical route modules with the
public import surface.  The key question is not how many modules exist, but
which assumptions remain theorem-facing.  In the current preferred API, the
answer is exactly the first-principles pair: chartwise smoothness and compact
closed support.
